% 両版共通のPDF内リンク。パッケージは他の設定の後で読み込む。
\usepackage[dvipdfmx,bookmarksnumbered,bookmarksopen,linktoc=all,hidelinks]{hyperref}
\usepackage{pxjahyper}

% 【例題】から【解答】へ、【解答】・【方針】・【解法】から【例題】へ。
% 第1引数は問題群ごとに固有のID。番号のない例題にも使える。
\newcommand{\bookblocklink}[4]{%
  \hypertarget{book:#1:#2}{\hyperlink{book:#1:#3}{\textbf{#4}}}%
}

% enumerate の format に指定し、既存の番号・インデントを保つ。
% 問題番号は解答へ、それ以外の番号は元の問題へリンクする。
% 各問題群の problem / answer / hint / solution で番号をそろえる。
\newcommand{\bookitemlink}[4]{%
  \hypertarget{book:#1:#2:\arabic{enumi}}{%
    \hyperlink{book:#1:#3:\arabic{enumi}}{#4}%
  }%
}
% 既存の問題番号を本文・学習経路から参照する。
% 第1引数は問題群ID、第2引数は番号、第3引数は表示する文字。
\newcommand{\bookproblemref}[3]{%
  \hyperlink{book:#1:problem:#2}{#3}%
}
